\chapter{Pricing-Library Architecture}\label{ch:pl:pricing-library-architecture}

A client asked on Friday for a price on a worst-of autocallable with memory coupons and a lookback on the knock-in: the capital is at risk
if the worst performance was ever below 60\% on an observation date, not only at maturity. The library had no class for it. By Monday it was
priced from eleven lines of \omterm{def:pl:pricing-library-architecture:script}{payoff script}, on the same \omterm{def:pl:pricing-library-architecture:market}{market objects}, model and Monte Carlo engine as every other trade in the library, with
Greeks, scenarios and batches for free. This chapter is about the architecture that makes that possible: a library in layers, whose engines
do not know which products they price, and a small language in which a product is data.

\section{The analytics library and its layers}

\begin{definition}[Analytics library]\label{def:pl:pricing-library-architecture:library}
An \emph{analytics library}\index{analytics library} is the firm's shared code for valuing and risk-managing instruments: \omterm{def:pl:pricing-library-architecture:market}{market objects},
models, numerical engines and instruments, separated so that each can be replaced or extended without the others, and called the same way by
pricing tools, risk systems and research.
\end{definition}

Book~5 built one (\texttt{firm.pricing}, chapter~28) and froze its interface for the rest of the series: instruments are descriptions of
cash flows; market data is an immutable snapshot; a model says how the market evolves; an engine computes a value for an instrument under a
model on a snapshot; and every Greek, scenario and batch is one repricing function applied to bumped snapshots. The layers are the
architecture: a new engine prices old instruments, a new instrument uses old engines, and a risk system (chapter~18, and Book~6's engine)
needs to know nothing but the repricing function.

\begin{omfigure}
\begin{tikzpicture}[font=\scriptsize,box/.style={draw,rounded corners=2pt,align=center,minimum height=0.8cm,inner sep=3pt},>=Stealth]
\node[box,fill=gray!10,minimum width=11cm] (api) at (0,2.2) {calls: \texttt{price}, \texttt{greeks}, \texttt{sensitivities},
\texttt{reprice}, \texttt{price\_batch} --- one repricing function on bumped snapshots};
\node[box,fill=omMeth!12,minimum width=2.6cm] (i) at (-4.7,0.9) {instruments\\ cash flows};
\node[box,fill=omMeth!12,minimum width=2.6cm] (m) at (-1.57,0.9) {market objects\\ curves, surfaces};
\node[box,fill=omMeth!12,minimum width=2.6cm] (mo) at (1.57,0.9) {models\\ Black--Scholes, Heston};
\node[box,fill=omMeth!12,minimum width=2.6cm] (e) at (4.7,0.9) {engines\\ analytic, PDE, MC};
\node[box,fill=omDef!15,minimum width=3.8cm] (s) at (-3.2,-0.6) {scripted instrument\\ script + schedules};
\node[box,fill=omDef!15,minimum width=3.8cm] (se) at (3.2,-0.6) {script engine\\ parse, schedule, evaluate};
\draw[->] (s.north) -- (s.north|-i.south); \draw[->] (se.north) -- (se.north|-e.south);
\draw[->] (se.west) -- (s.east);
\foreach \x in {i,m,mo,e} {\draw[->] (\x.north) -- (\x.north|-api.south);}
\node[font=\scriptsize\itshape] at (0,-1.55) {registered with \texttt{firm.pricing.register}: priced, bumped and batched like any other};
\end{tikzpicture}

\omcaption{The library's layers and the script engine's place in them. A scripted instrument is an instrument; the script engine is an
engine, registered for the Black--Scholes model; everything above --- Greeks, scenarios, batches --- applies unchanged.}\label{fig:pl:pricing-library-architecture:layers}
\end{omfigure}

\section{Instruments, market objects, models and engines}

\begin{definition}[Market object]\label{def:pl:pricing-library-architecture:market}
A \emph{market object}\index{market object} is a library object that represents one piece of market state --- a quote, a discount curve, a
volatility surface, a correlation --- and answers questions about it (a discount factor, an implied volatility), so that engines read the
market through a fixed interface whatever its construction.
\end{definition}

In \texttt{firm.pricing} a snapshot holds the \omterm{def:pl:pricing-library-architecture:market}{market objects} and a bump makes a new snapshot with one object replaced; the Monte Carlo engine
reads forwards, variances and correlations from it and generates paths from a seed fixed per instrument (common random numbers, Book~5), so
that a bumped price differs from the base price by the bump alone. The script engine of this chapter reuses those paths: it is the Monte Carlo
engine with the payoff taken from a script instead of from a class.

\section{Payoff scripting}

\begin{definition}[Payoff scripting language, payoff script]\label{def:pl:pricing-library-architecture:script}
A \emph{payoff scripting language}\index{payoff scripting language} is a small programming language in which a product's cash flows are
written as events on observation dates --- conditions on the path, payments, state carried from date to date --- and evaluated by a generic
engine on simulated paths. A \emph{payoff script}\index{payoff script} is one product written in it: data that the library prices, not code
that it must be extended with.
\end{definition}

The chapter's language (\texttt{firm.payoffdsl}) has what the chapter's products need and nothing more: state variables, blocks run at the
dates of named schedules, assignments, \texttt{if}/\texttt{else}, \texttt{pay} and \texttt{stop}, arithmetic and comparisons, and
\texttt{min}, \texttt{max} and \texttt{abs}. At each date a block sees \texttt{S} (the first underlying's spot), \texttt{perf} (the worst
performance) and \texttt{t}. A tokenizer turns the text into tokens, a recursive-descent parser turns them into a syntax tree with an error
position for anything it cannot read, and the schedule is extracted from the tree: every block's dates, in date order, in script order on the
same date.

\begin{omfigure}
\begin{tikzpicture}[font=\scriptsize,box/.style={draw,rounded corners=2pt,align=center,minimum height=0.9cm,minimum width=1.9cm,
  fill=omDef!10},>=Stealth]
\node[box] (a) at (0,0) {script text};
\node[box] (b) at (2.4,0) {tokens\\ with positions};
\node[box] (c) at (4.8,0) {syntax tree};
\node[box] (d) at (7.2,0) {event schedule\\ (date, block)};
\node[box,fill=omMeth!12] (e1) at (10.4,0.6) {interpreter\\ per path};
\node[box,fill=omMeth!12] (e2) at (10.4,-0.6) {compiled\\ arrays, masks};
\draw[->] (a) -- (b); \draw[->] (b) -- (c); \draw[->] (c) -- (d); \draw[->] (d) -- (e1.west); \draw[->] (d) -- (e2.west);
\node[font=\scriptsize\itshape,omThm] at (3.4,-0.95) {errors with line and column};
\node[font=\scriptsize\itshape] at (7.2,0.95) {library paths on these dates};
\end{tikzpicture}

\omcaption{From a script to a price: the tokenizer and the parser build the tree, the schedule is read from its blocks, and either evaluator
runs the tree on the library's paths for the schedule's dates.}\label{fig:pl:pricing-library-architecture:pipeline}
\end{omfigure}

\omcode{code/platforms/19-pricing-library-architecture/python/pl_payoffdsl.py}{36}{44}{The worst-of Phoenix autocallable with memory
coupons as a payoff script: coupons on every observation date, an autocall on all but the last, the capital at risk at
maturity.}\label{lst:pl:pricing-library-architecture:autocall}

\omcode{code/firm/payoffdsl/firm_payoffdsl.py}{176}{199}{The parser's statements: \texttt{if}/\texttt{else}, \texttt{pay}, \texttt{stop}
and assignment; every token keeps its line and column for error messages.}\label{lst:pl:pricing-library-architecture:parser}

The client's product changes three lines: a state variable \texttt{lowest}, updated with \texttt{lowest = min(lowest, perf)} at every
observation, and the maturity test on \texttt{lowest} instead of \texttt{perf}. It prices at 94.25 per 100, 1.63 below the plain
autocallable's 95.88: the lookback knock-in triggers on paths that dip below 60\% and recover.

\section{Two ways to evaluate a script}

The simplest evaluator walks the syntax tree once per path and per event, with scalar values: an interpreter. It is easy to trust and slow.
The compiled evaluator walks the tree once per event with arrays over all paths: every statement runs under a mask of the paths it applies to,
an \texttt{if} splits the mask, a \texttt{pay} adds under it, a \texttt{stop} removes paths from the living set
(\cref{lst:pl:pricing-library-architecture:compile}). The two give the same numbers to the last digit (the tests compare them); the compiled one
is 19 to 34 times faster (\cref{fig:pl:pricing-library-architecture:speed}).

\omcode{code/firm/payoffdsl/firm_payoffdsl.py}{357}{388}{The compiled evaluation: masks instead of branches, arrays instead of
loops.}\label{lst:pl:pricing-library-architecture:compile}

\begin{omfigure}
\begin{tikzpicture}
\begin{loglogaxis}[width=11cm,height=5.6cm,xlabel={paths (log scale)},ylabel={seconds per price (log scale)},grid=major,grid style={gray!20},
  tick label style={font=\small},label style={font=\small},table/col sep=comma,legend style={font=\footnotesize,at={(0.5,1.03)},
  anchor=south,legend columns=2}]
\addplot[color=omThm,very thick,mark=*] table[x=paths,y=interpreted_s]{figdata/platforms/19-pricing-library-architecture/measured_speed.csv};
\addlegendentry{interpreted (per path)}
\addplot[color=omDef,very thick,mark=square*] table[x=paths,y=compiled_s]{figdata/platforms/19-pricing-library-architecture/measured_speed.csv};
\addlegendentry{compiled (arrays over paths)}
\end{loglogaxis}
\end{tikzpicture}

\omcaption{Time to price the autocallable script by the interpreter and by the compiled evaluator, path generation included, best of three,
one thread. The compiled version is 18.8 times faster at 1\,000 paths, 34.3 times at 10\,000 and 33.4 times at 30\,000. Measured on a laptop (Intel Core Ultra 7
155H) under WSL2, machine otherwise idle. Data: \texttt{bench\_payoffdsl.py}.}\label{fig:pl:pricing-library-architecture:speed}
\end{omfigure}

The script engine is registered in the library like any engine (\cref{lst:pl:pricing-library-architecture:engine}): it parses the script,
extracts the schedule, asks the Monte Carlo engine for paths on the schedule's dates with the instrument's seed, and evaluates. Because the
seed depends only on the instrument's identifier, a script with the same identifier as a library instrument sees the same paths.

\omcode{code/firm/payoffdsl/firm_payoffdsl.py}{403}{429}{The script engine: schedule, the library's paths, discounting on the instrument's
curve, and the evaluator.}\label{lst:pl:pricing-library-architecture:engine}

\begin{table}[ht]
\centering\small
\begin{tabular}{lrrrrl}
\toprule
product & script & s.e. & library & difference (s.e.) & library engine \\
\midrule
European call & 11.3953 & 0.0570 & 11.3485 & 0.82 & analytic \\
Asian call (12 monthly fixings) & 6.8012 & 0.0327 & 6.8012 & 0.00 & Monte Carlo \\
daily down-and-out call & 11.0373 & 0.0568 & 11.0226 & 0.26 & closed form (BGK) \\
worst-of autocallable & 95.8750 & 0.0717 & 95.8750 & 0.00 & Monte Carlo \\
autocallable, lookback knock-in & 94.2474 & 0.0726 & --- & --- & none \\
\bottomrule
\end{tabular}
\caption{Scripts against the library's own engines, 100\,000 paths. Against an analytic or closed-form engine the difference is Monte Carlo
error, under one standard error; against the library's Monte Carlo engine on common random numbers it is zero, because the paths and the
payoff arithmetic are the same.}\label{tab:pl:pricing-library-architecture:validation}
\end{table}

\begin{example}[Priced by Monday]\label{ex:pl:pricing-library-architecture:monday}
The scripted autocallable prices at 95.8750 per 100 and the library's own autocallable engine at 95.8750: the difference is 0.00 standard
errors, because both evaluate the same payoff on the same paths. Their deltas to the first underlying are equal too, 0.1518 per unit of spot
per 100 of notional, since the library's bump-and-reprice Greeks work on any engine. The client's lookback variant, which the library cannot
express, prices at 94.2474 ($\pm0.0726$) with the same machinery. The compiled evaluator makes the script engine as fast as a hand-written
array engine; the interpreter, 22 to 37 times slower, remains the reference that checks it.
\end{example}

\section{Lazy recalculation and observers}

\begin{definition}[Observer pattern]\label{def:pl:pricing-library-architecture:observer}
The \emph{observer pattern}\index{observer pattern} lets an object (the observable) keep a list of dependants (observers) and notify them
when it changes, without knowing what they are; a curve observes its quotes, a price observes its curve.
\end{definition}

\begin{definition}[Lazy recalculation]\label{def:pl:pricing-library-architecture:lazy}
\emph{Lazy recalculation}\index{lazy recalculation} marks a result stale when a notification arrives and recomputes it only when the result
is next asked for, so that a burst of market changes costs one recalculation per reader instead of one per change.
\end{definition}

QuantLib's \texttt{LazyObject} is the best-known implementation: a ``framework for calculation on demand and result caching'' that is both
an observer and an observable. The chapter's \texttt{LazyPrice} is the same idea in twenty lines
(\cref{lst:pl:pricing-library-architecture:lazy}): a European call observing a rate quote is computed once (11.35), the quote changes three
times, and the next read recomputes once (11.59) --- two calculations for four states of the market. The snapshot design of
\texttt{firm.pricing} is the other answer to the same problem: nothing is mutable, so nothing is stale; a new snapshot is a new question.
Libraries serving interactive tools favour observers; libraries serving batch risk favour snapshots.

\omcode{code/firm/payoffdsl/firm_payoffdsl.py}{469}{490}{A lazy price: stale on notification, recomputed on demand, and itself observable
by what depends on it.}\label{lst:pl:pricing-library-architecture:lazy}

\section{Bindings to other languages}

A library written once is called from many places: Python research, a C++ pricing service, a spreadsheet on a trader's desk. Chapter~9
measured what a binding costs per call; for a library the lesson is to bind at the level of whole requests --- a batch of trades, a snapshot
--- never at the level of an operation. Book~5's library has C++20 and Rust twins of its core, tested against the Python reference on the
same inputs; a \omterm{def:pl:pricing-library-architecture:script}{payoff script} is portable by construction, since it is text that any implementation of the language can price, which is why
scripts travel between front-office tools, risk systems and the canonical trade model of chapter~21 more easily than classes do.

\section{Tutorial: a language, an engine, five products}\label{tut:pl:pricing-library-architecture:tut}

\begin{tutorial}
\textbf{Goal.} Write products as scripts, price them through the library, check them against the library's engines, and measure the two
evaluators. \textbf{End state:} \cref{tab:pl:pricing-library-architecture:validation} and \cref{fig:pl:pricing-library-architecture:speed}.
\begin{tutsteps}
\item \textbf{Parse}: \texttt{firm\_payoffdsl.parse} on each script; break one and read the error position.
\item \textbf{Price}: \texttt{firm.pricing.price} on a \texttt{ScriptedInstrument}; \texttt{greeks} on it.
\item \textbf{Validate}: \texttt{pl\_payoffdsl.validation()}.
\item \textbf{Measure}: \texttt{bench\_payoffdsl.py}.
\item \textbf{Observe}: \texttt{lazy\_demo()}.
\end{tutsteps}
\textbf{What to change next.} Add a \texttt{worst} function over named underlyings and a daily lookback between observation dates; compile a
script to Python source instead of closures and measure again.
\end{tutorial}

\section{Build: the payoff language}\label{bld:pl:pricing-library-architecture:payoffdsl}

\begin{build}
\bfield{Purpose} Products the library has no class for, priced, bumped and batched like those it has, from scripts that can be stored,
reviewed and moved between systems.
\bfield{Interface} \texttt{parse}, \texttt{schedule}, \texttt{Interpreter}, \texttt{compile\_program}, \texttt{ScriptedInstrument},
\texttt{ScriptEngine} (registered for Black--Scholes), \texttt{ScriptError}; \texttt{Quote}, \texttt{LazyPrice}.
\bfield{Rules} The library's paths and seeds, never the engine's own; discounting on the instrument's curve; the interpreter and the compiled
evaluator agree on every path; errors name a line and a column.
\bfield{Acceptance tests} \texttt{code/firm/payoffdsl/tests/}: parsing and error positions; the event schedule's order; the interpreter equal
to the compiled evaluator on a path-dependent script; a scripted European equal to the library's Monte Carlo price and close to its analytic
one, with its delta; lazy prices and chained observers.
\bfield{Stretch} Several underlyings by name; adjoint (AAD) Greeks through the compiled script (Book~4, chapter~28); a script library with
versioned, reviewed products.
\end{build}

\begin{omsources}
\item QuantLib reference documentation, \texttt{LazyObject}.
\item One Quant Book~5, chapters 4, 18 and 28 (Greeks, autocallables, the library); One Quant Book~4, chapter~26 (Monte Carlo).
\end{omsources}

\section{Exercises}

\begin{exercise}[$\star$]\label{exo:pl:pricing-library-architecture:1}
Write a digital call paying 10 at expiry if the spot is above 105 as a \omterm{def:pl:pricing-library-architecture:script}{payoff script}.
\end{exercise}

\begin{exercise}[$\star$]\label{exo:pl:pricing-library-architecture:2}
Why is the scripted autocallable's price exactly the library's, while the scripted European differs from the analytic price by 0.82
standard errors?
\end{exercise}

\begin{exercise}[$\star$]\label{exo:pl:pricing-library-architecture:3}
In the autocallable script, why must the coupon block come before the autocall block on the same date?
\end{exercise}

\begin{exercise}[$\star\star$]\label{exo:pl:pricing-library-architecture:4}
How does the compiled evaluator handle \texttt{stop}, and what would go wrong if it simply skipped the rest of the block?
\end{exercise}

\begin{exercise}[$\star\star$]\label{exo:pl:pricing-library-architecture:5}
Why does the speed-up of the compiled evaluator grow with the number of paths and then level off?
\end{exercise}

\begin{exercise}[$\star\star$]\label{exo:pl:pricing-library-architecture:6}
Observers or snapshots: which would you choose for a trader's pricing screen, and which for the overnight risk run?
\end{exercise}

\begin{exercise}[$\star\star\star$]\label{exo:pl:pricing-library-architecture:7}
\emph{Coding.} Price the lookback autocallable with the knock-in level at 50\% and at 70\%. How does the price move, and what does that say
about the product's sensitivity to the barrier?
\end{exercise}

\begin{exercise}[$\star\star\star$]\label{exo:pl:pricing-library-architecture:8}
\emph{Find the flaw.} ``Our scripting engine draws its own random numbers, so script prices are independent of the library and a good
cross-check.''
\end{exercise}

\section{Problem: Priced by Monday}

\begin{problem}[{Weekend problem --- a product the library has no class for}]\label{pb:pl:pricing-library-architecture:1}
\texttt{firm.pricing}, \texttt{firm.payoffdsl} and the chapter's five products.

\medskip\noindent\textbf{Part I --- The library.}
\begin{enumerate}
\item What are the library's four layers, and what connects them?
\item What does a \omterm{def:pl:pricing-library-architecture:market}{market object} answer, and why an immutable snapshot?
\item Why do bumped prices differ by the bump alone?
\item What must an engine provide to be registered?
\item Why does a risk system need nothing but the repricing function?
\end{enumerate}

\medskip\noindent\textbf{Part II --- The language.}
\begin{enumerate}[resume]
\item What can a script say, and what does each block see?
\item How is the event schedule extracted?
\item How do the interpreter and the compiled evaluator differ?
\item How does the client's product differ from the plain autocallable in the script?
\item What happens to a syntax error?
\end{enumerate}

\medskip\noindent\textbf{Part III --- The numbers.}
\begin{enumerate}[resume]
\item Give the script and library prices of the five products and their differences in standard errors.
\item Why are two differences exactly zero?
\item What are the autocallable's deltas by script and by library?
\item What is the compiled evaluator's speed-up?
\item What does the lazy price's demonstration count?
\end{enumerate}

\medskip\noindent\textbf{Part IV --- The verdict.}
\begin{enumerate}[resume]
\item State the \emph{named result}: the scripted autocallable's price and delta against the library's autocallable engine on common random
numbers, their difference in standard errors, and the speed-up of the compiled script over the interpreted one.
\item What would you require before a scripted product is traded?
\item When should a product get its own class and engine instead of a script?
\item How would you store and version scripts?
\item In one sentence: what does a payoff language buy a pricing library?
\end{enumerate}
\end{problem}

\section{Interview questions}

\begin{interviewq}[$\star$ \iqroles{developer}]\label{iq:pl:pricing-library-architecture:1}
Why separate instruments, models and engines in a pricing library?
\end{interviewq}

\begin{interviewq}[$\star\star$ \iqroles{developer, researcher}]\label{iq:pl:pricing-library-architecture:2}
How would you price a new exotic by Monday without writing a new engine?
\end{interviewq}

\begin{interviewq}[$\star\star$ \iqroles{developer}]\label{iq:pl:pricing-library-architecture:3}
How do you make a Monte Carlo payoff evaluator fast in Python?
\end{interviewq}

\begin{interviewq}[$\star\star$ \iqroles{developer}]\label{iq:pl:pricing-library-architecture:4}
What is \omterm{def:pl:dataframes:eager}{lazy evaluation} in a pricing library, and when is it a bad idea?
\end{interviewq}

\begin{interviewq}[$\star\star\star$ \iqroles{developer}]\label{iq:pl:pricing-library-architecture:5}
Design a \omterm{def:pl:pricing-library-architecture:script}{payoff scripting language} for a derivatives desk.
\end{interviewq}
